\chapter{La partícula como sección material persistente: historia, representación, banda, composición y criterio de identidad}
\label{chap:particula-seccion-material-persistente}

\section{Definición de partícula como clase compatible de secciones materiales persistentes bajo refinamiento}
Una partícula es una clase compatible de secciones materiales persistentes de
la extensión genealógica, dotada de historia observable, representación,
orientación y criterio de identidad bajo refinamiento. La ruta es su historia,
no la partícula misma; la celda es una coordenada del atlas, no un cuerpo que
se desplaza; la masa es un observable espectral, no su definición ontológica.

\section{Clasificador de especies materiales por fibra, representación, carga, espín, sabor y generación}
La especie se tipa mediante carga, espín, color, sabor, generación,
representación y orientación. La carga es un carácter orientado; el espín es
un dato de representación; el color es una fibra ternaria residual; sabor y
generación pertenecen a la multisección y a su profundidad. Ninguno de estos
datos se reduce al valor de masa.

\section{Antipartícula, banda y virtualidad}
La involución \(C_M=-\operatorname{rev}\) induce la antisección y la
conjugación de hojas. Partícula y antipartícula pueden ser dos secciones de una
misma cubierta de Möbius cuando se fijan el lazo, la holonomía y la
representación correspondientes. El retorno extendido de 216 pasos no se
descompone en dos trayectorias independientes de 108. Una partícula virtual es
una sección intermedia no asintótica con horizonte finito de prolongación.

\section{Persistencia de trayectoria, multisecciones y conjugación de Möbius}

La persistencia se formula como sección de un sistema inverso; la especie
interna como cociente orbital de una multisección; y la conjugación como
anti-sección con holonomía declarada. El desarrollo completo conserva la
estratificación extensión--ruta--registro--firma, el lector temporal y la
distinción entre las dos apariciones de \(2\pi/4\pi\).

\begingroup
\let\chapter\section
\let\section\subsection
\let\subsection\subsubsection
\let\subsubsection\paragraph
\let\clearpage\relax
\let\cleardoublepage\relax
\section{Persistencia de ruta, multisecciones de especie y conjugaci\'on de M\"obius}
\label{sec:persistencia-ruta-mobius-familias}

Esta secci\'on precisa una distinci\'on que, si se formula de manera demasiado
abreviada, invierte la l\'ogica de la construcci\'on. La ruta no es un dato que
una part\'icula escoge una vez conocida su masa. Tampoco es un dibujo auxiliar
sin contenido causal. Es la historia observable y orientada de una extensi\'on
superviviente. La extensi\'on conserva, adem\'as, la profundidad compatible, el
registro, los acarreos y la memoria que no caben en un corte finito. En este
sentido preciso,
\[
 \boxed{\text{una part\'icula HMT es la persistencia compatible de una ruta
 enriquecida}.}
\]
La frase y la definici\'on por extensiones expresan el mismo objeto a dos
niveles: la ruta es su historia visible; la extensi\'on es la torre completa de
esa historia con todas sus prolongaciones compatibles.

\subsection{Persistencia como secci\'on de un sistema inverso}

Sea \((\mathcal S_m,\rho_{m+1,m})_{m\geq0}\) el sistema de estados que
sobreviven hasta profundidad \(m\), y sea
\[
 x=(x_m)_{m\geq0}\in\Omega_{\rm HMT}^{\rm surv}
 =\varprojlim_m\mathcal S_m.
\]
Las sombras de ruta satisfacen el cuadrado de compatibilidad
\[
 \operatorname{tr}_{m+1,m}\!\left(\operatorname{sh}_{m+1}(x_{m+1})\right)
 =\operatorname{sh}_{m}\!\left(\rho_{m+1,m}(x_{m+1})\right)
 =\operatorname{sh}_{m}(x_m).
\]
Por tanto
\[
 \gamma_x=\bigl(\gamma_x^{[m]}\bigr)_{m\geq0},
 \qquad \gamma_x^{[m]}=\operatorname{sh}_m(x_m),
\]
es una historia compatible, no una secuencia ajustada retrospectivamente.

Sobre cada corte \(\gamma_x^{[m]}\) se dispone del fibrado discreto
\(\mathcal B_{\gamma_x^{[m]}}\). Una secci\'on persistente es una familia
\[
 s=(s_m)_{m\geq0},
 \qquad
 s_m\in\Gamma\!\left(\mathcal B_{\gamma_x^{[m]}}\right),
 \qquad
 \operatorname{res}_{m+1,m}(s_{m+1})=s_m.
\]
Esta ecuaci\'on es el contenido matem\'atico de la persistencia: cada lectura
finita se prolonga sin contradecir las lecturas anteriores. La tupla de
part\'icula del \cref{chap:particula} a\~nade ocupaci\'on,
representaci\'on y propagaci\'on, pero no altera esta procedencia.

\subsection{Especie interna como cociente orbital de la multisección}

En el atlas non\'adico, cada celda orientada \(c=(r,c')\in\mathcal C_{81}\)
porta una familia de ruta
\[
 \mathcal R(c)=
 (n_A,s_C,k_{\Delta_4},\nu_{120},\nu_{270},\tau_{12}).
\]
La imagen de \(\mathcal R\) contiene cincuenta y seis valores. Si se fija la
carta y su orden, el rango \(\kappa\in\{0,1,2,3,4\}\) dentro de cada fibra
produce el levantamiento inyectivo
\[
 c\longmapsto \bigl(\mathcal R(c),\kappa(c)\bigr).
\]
El alfabeto de cinco estados es m\'inimo, pues existen dos fibras de
cardinal cinco. Esta selecci\'on es interna, prospectiva y ciega a masas y a
nombres de part\'iculas.

El censo completo que sostiene esta afirmaci\'on es
\[
 81=45_{\rm leptones}+36_{\rm quarks}
   =25_{\ell^-}+20_{\nu}+20_{d\text{-tipo}}+16_{u\text{-tipo}}.
\]
Las cincuenta y seis fibras de ruta poseen el histograma
\[
 41\times1+10\times2+2\times3+1\times4+2\times5=81.
\]
Por tanto, \(56\) cuenta firmas de ruta, mientras \(81\) cuenta celdas
orientadas; la memoria \(\kappa\) distingue las celdas que una misma firma no
separa.

Para clasificar especies internas no se invierte este mapa. Se toma el
cociente
\[
 q:\mathcal C_{81}\longrightarrow\Sigma_{13},
 \qquad
 q(c)=\bigl(\operatorname{familia}(c),G(c)\bigr),
\]
cuya imagen comprende trece tipos. La salida asociada a un tipo \(y\) es la
multisecci\'on finita del atlas
\[
 \boxed{
 \mathfrak S_y^{\rm atlas}=
 \{(\mathcal R(c),\kappa(c)):q(c)=y\}.}
\]
Sus cardinales son
\[
 1,8,16;\qquad 2,4,6,8;\qquad 2,4,6,8;\qquad4,12,
\]
para leptones cargados, neutrinos y los dos tipos de quark,
respectivamente. La especie es el tipo com\'un de esa fibra; una part\'icula
concreta es una secci\'on persistente realizada dentro de ella. Exigir que el
nombre de una especie elija por s\'i solo una celda destruye precisamente los
grados de orientaci\'on y conjugaci\'on que la fibra conserva.

La elevaci\'on a extensiones es un objeto tipado diferente. Si
\(\pi_{81}:\Omega_{\rm HMT}^{\rm surv}\to\mathcal C_{81}\) denota el lector
celular de la ruta enriquecida declarada, entonces
\[
 \widetilde{\mathfrak S}_y
 =\{x\in\Omega_{\rm HMT}^{\rm surv}:q(\pi_{81}x)=y\}
\]
es la multisecci\'on enriquecida correspondiente. El certificado finito
demuestra el censo y la obstrucci\'on en
\(\mathfrak S_y^{\rm atlas}\); su elevaci\'on es relativa al lector
\(\pi_{81}\), no una identificaci\'on autom\'atica entre una celda y una
historia profunda.

\begin{theorem}[Obstrucci\'on al selector puntual y excepci\'on electr\'onica]
Sea \(\rho_{\rm c}(r,c')=(10-r,10-c')\) la inversi\'on central. Las trece fibras de
\(q\) son \(\rho_{\rm c}\)-estables. La fibra de lept\'on cargado de nivel cero posee
el \'unico punto fijo, \((5,5)\). Ninguna de las otras doce fibras admite
un selector puntual \(\rho_{\rm c}\)-equivariante si \(\rho_{\rm c}\) act\'ua trivialmente
sobre el tipo grueso.
\end{theorem}
\begin{proof}
La ecuaci\'on \(\rho_{\rm c}(r,c')=(r,c')\) equivale a \(r=c'=5\). El censo del atlas
sit\'ua esa celda en la fibra \((\ell^-,G0)\), que tiene cardinal uno. Sea
ahora \(y\neq(\ell^-,G0)\) y supongamos que existe una secci\'on puntual
equivariante \(a:\Sigma_{13}\to\mathcal C_{81}\). Como \(\rho_{\rm c} y=y\), la
equivarianza exigir\'ia
\[
 a(y)=a(\rho_{\rm c} y)=\rho_{\rm c} a(y).
\]
Luego \(a(y)\) tendr\'ia que ser un punto fijo de \(\rho_{\rm c}\), pero el \'unico
punto fijo pertenece a la fibra electr\'onica. Esto es imposible. Por tanto,
en las otras doce fibras la salida natural es una \(\rho_{\rm c}\)-\'orbita o una
multisecci\'on; elegir un representante requiere orientaci\'on o
representaci\'on adicional.
\end{proof}

El electr\'on es excepcional en tres sentidos compatibles: \((5,5)\) es el
centro geom\'etrico \'unico, la \'unica celda de nivel cero y el centro respecto
del cual se forma el ret\'iculo crudo de diferencias
\[
 \Lambda_{\rm raw}^{\Delta}
 =\bigl\langle A_{\rm raw}(c)-A_{\rm raw}(5,5):
 c\in\mathcal C_{81}\bigr\rangle_{\mathbb Z}.
\]
La firma cruda de esa celda no es nula:
\[
 A_{\rm raw}(5,5)=(24,0,12,0,-12),
 \qquad \tau(5,5)=\texttt{+++---+++---}.
\]
El origen \(v_e=(0,0,0,0,0)\) del ret\'iculo de acci\'on es una convenci\'on
relativa al electr\'on y no la firma cruda de la celda central.

El invariante algebraico electr\'onico se eval\'ua mediante
\[
 \mathfrak e_0=
 \frac{\sqrt3}{4}
 \left(e^{\varphi/\pi^2}-22\alpha_{\rm HMT}^{3}\right)
 (1+15\Delta_4),
 \qquad
 \Delta_4=
 \frac{\pi-e\log\pi}{270}\left(1+\frac{\pi}{729}\right),
\]
\[
 \mathfrak e_0
 =0.510998922488066072509870877191349\ldots.
\]
La expresi\'on es adimensional; la carta
\(\varsigma_{u_E}:y\mapsto y\,u_E\) introduce despu\'es una unidad de energ\'ia.
El centro \((5,5)\) y su unicidad son resultados finitos exactos. La selecci\'on
de los coeficientes \(22\) y \(15=3\cdot5\) se deduce, con la estructura de
partida expl\'icita, mediante el complemento
de las cinco regiones de \(\pi\) en \(K_e=\mathbb Z/9\mathbb Z\times\mathbb F_3\):
\(-22=-|K_e\setminus\iota(R_\pi)|\) y
\(+15=|R_\pi\times\mathbb F_3|\). La firma cruda central conserva un tipo
distinto y act\'ua como control independiente de la selecci\'on basal.

La celda \((5,5)\) no debe confundirse con la firma energ\'etica
\(555555\): esta \'ultima pertenece al representante transversal de una de
las cinco regiones de \(\pi\). Son objetos distintos que comparten el d\'igito
central.

En el sector de quarks, el car\'acter residual
\[
 \operatorname{col}(r,c')=r-c'\pmod3
\]
refina la multisecci\'on en tres clases de doce celdas y satisface
\(\operatorname{col}(\rho_{\rm c} c)=-\operatorname{col}(c)\). Esto construye el
tr\'itico interno de color. La representaci\'on f\'isica de \(SU(3)\) es un
lector posterior; no hace falta para demostrar el reparto finito.

\subsection{Estratificación de extensión, ruta, registro y firma}

Los mapas \(q\), \(\mathcal R\), la proyecci\'on celular y el lector de
realizaci\'on distinguen cuatro clasificaciones relacionadas y tipadas:
\begin{enumerate}[label=(\roman*)]
 \item las trece fibras internas familia--nivel de \(q\);
 \item las cincuenta y seis familias de ruta de \(\mathcal R\);
 \item las ochenta y una posiciones y sus trescientas veinticuatro
       orientaciones;
 \item las realizaciones f\'isicas individualizadas por representaci\'on y codominio.
\end{enumerate}
Las cuatro capas se censan antes del contraste metrol\'ogico. Para cada clase,
la comparación requiere la tupla
\[
(\text{posici\'on},\text{fibra},\text{rutas},\text{lectores},
\text{registro},\text{firma},\text{car\'acter},\text{torres},
\text{carta dimensional}).
\]
Una reconstrucci\'on sobre un subconjunto sólo demuestra ese subcaso y no
reemplaza ninguno de los componentes de la tupla.
Dados \(A\), \(C^*\) y los canales \(q_\pm\), la torre global pertenece al
semigrupo infinito
\[
 s_n=q_-^n-q_+^n,
 \qquad
 n\in\langle90,120\rangle
 =\{90,120\}\cup\{30r:r\geq6\}.
\]
Las alturas \(90,120,180,210,240,270,360,420\) son representantes finitos del
semigrupo; la recurrencia anterior prolonga la torre completa.
y el car\'acter adimensional de una firma es
\[
 \chi_{\rm mass}(v)=\exp\!\left(
 n_AA_{\rm rad}+\frac{s_C}{6}C^*_{\rm rad}
 +k\Delta_4+\nu_{120}s_{120}+\nu_{270}(135\Delta_4^2)
 \right).
\]
Los momentos son globales: las especies los ponderan mediante su ruta, su
registro y su operador de fibra. Los lectores \(K_D\) y \(K_\Omega\)
determinan el refinamiento bidireccional antes del contraste, mientras la
tabla integral de especies publica por separado la genealogía interna, la
carta dimensional y la comparación externa.

\subsection{Conjugaci\'on, anti-secci\'on y banda de M\"obius}

Sobre una palabra dodecaf\'asica se define la involuci\'on exacta
\[
 C_M(\tau)_m=-\tau_{11-m},
 \qquad C_M^2=I.
\]
Los funcionales lineales con pesos sim\'etricos son impares bajo \(C_M\),
mientras las correlaciones cuadr\'aticas compatibles con la reversi\'on son
pares. De este modo, orientaci\'on y carga lineal viven en el sector impar,
mientras la memoria cuadr\'atica vive en el sector par. Esto no demuestra que
el car\'acter exponencial completo de masa sea par: si
\(\chi(\tau)=e^{\ell(\tau)}\) y \(\ell\) es impar, entonces
\(\chi(C_M\tau)=\chi(\tau)^{-1}\).

La involuci\'on anterior act\'ua exactamente sobre palabras. Para hablar de
anti-secci\'on se fija adem\'as un levantamiento \(\widetilde C_M\) al fibrado de
secciones que preserve admisibilidad, registro y truncamientos:
\[
 \operatorname{res}_{m+1,m}\widetilde C_M
 =\widetilde C_M\operatorname{res}_{m+1,m},
 \qquad
 \operatorname{sh}(\widetilde C_Mx)=C_M\operatorname{sh}(x).
\]
S\'olo bajo ese levantamiento la anti-secci\'on de \(s\) es
\(\widetilde C_Ms\) sobre la ruta invertida.

La palabra ``M\"obius'' tiene aqu\'i una formulaci\'on relativa precisa. Sea
\(\ell_{\rm ret}\) un lazo de retorno de la ruta enriquecida que genera
\[
 \langle\ell_{\rm ret}\rangle\simeq\pi_1(S^1)\simeq\mathbb Z.
\]
Este grupo de retorno no es la fase finita \(C_9\). Supongamos que su
levantamiento al
registro posee un car\'acter de signo
\[
 \varepsilon_x:\langle\ell_{\rm ret}\rangle\longrightarrow\{\pm1\},
 \qquad \varepsilon_x(\ell_{\rm ret})=-1.
\]
El fibrado de intervalos asociado es
\[
 \mathcal M_x=
 ([0,1]\times[-1,1])/(0,u)\sim(1,-u).
\]
Su primera clase de Stiefel--Whitney es no trivial; por tanto
\(\mathcal M_x\) es una banda de M\"obius. Una vuelta intercambia las dos
hojas locales y dos vueltas restauran la orientaci\'on. Esto no obliga a que
el registro completo vuelva al mismo estado: retorno de orientaci\'on y retorno
de memoria siguen siendo tipos diferentes.

El teorema es relativo al levantamiento, al lazo de retorno y al car\'acter
\(\varepsilon_x\), todos ellos declarados. No nace de observar dos signos.
La lectura de hojas requiere, adem\'as, distinguir la graduaci\'on partida \(R\)
de la conjugaci\'on \(C\) e imponer
\[
 CRC^{-1}=-R,
 \qquad CT(C^*)C^{-1}=T(-C^*).
\]
La lectura part\'icula--antipart\'icula exige finalmente que la representaci\'on
f\'isica invierta los portadores impares, preserve los portadores cuadr\'aticos
pares y entrelace la evoluci\'on y el car\'acter completo de la forma declarada
por la representaci\'on. Bajo esas premisas, part\'icula y antipart\'icula son
las dos orientaciones conjugadas de una misma banda de persistencia.

\begin{theorem}[Descomposici\'on par--impar sobre una cubierta doble]
\label{pf84:A10-md-mobius}
Sea \(p:\widetilde X\to X\) una cubierta principal doble con involuci\'on
\(\iota\). Para toda \(f\in C^0(\widetilde X,\mathbb R)\), los proyectores
\[
 f^+=\frac12(f+\iota^*f),\qquad
 f^-=\frac12(f-\iota^*f)
\]
dan la descomposici\'on \(f=f^++f^-\), con
\[
 f^+\circ\iota=f^+,\qquad f^-\circ\iota=-f^-.
\]
La parte par desciende a una funci\'on sobre \(X\); la parte impar es una
secci\'on de la l\'inea de signo asociada. Para la cubierta de orientaci\'on
de la banda de M\"obius, \(w_1\ne0\) obstruye una secci\'on global continua
y no nula de esa l\'inea.
\end{theorem}

\begin{proof}
Las identidades de los dos proyectores prueban la suma y las paridades. Una
funci\'on invariante es constante en las fibras de \(p\), por lo que
desciende; una funci\'on antiinvariante cumple la ley de transici\'on de la
representaci\'on signo. En la banda de M\"obius la l\'inea asociada es no
trivial y su primera clase de Stiefel--Whitney es no nula, de modo que toda
secci\'on continua global tiene alg\'un cero. La conclusi\'on excluye una
secci\'on global continua \emph{sin ceros}; no niega la existencia de
secciones con ceros.
\end{proof}

Este lema cl\'asico tipa el sector par--impar de la cubierta; no identifica
por s\'i solo la reversi\'on de palabra, la holonom\'ia de retorno, la
graduaci\'on partida, la conjugaci\'on de hojas ni la paridad de una firma de
masa. Esas identificaciones requieren entrelazadores propios.

\subsection{TRIT algebraico y lector temporal}

La realizaci\'on cuadr\'atica declarada asocia al signo visible del TRIT un tipo
algebraico. Para
\(\bar\tau\in\{+1,0,-1\}\), sea
\[
 \mathcal A_{\bar\tau}=
 \mathbb R[J]/(J^2+\bar\tau).
\]
Entonces
\[
 \begin{array}{c|c|c}
 \bar\tau&J^2&\text{r\'egimen}\ \\ \hline
 +1&-1&\text{el\'iptico/complejo}\ \\
 0&0&\text{parab\'olico/nilpotente}\ \\
 -1&+1&\text{hiperb\'olico/partido}.
 \end{array}
\]
Las exponenciales correspondientes son
\[
 e^{tJ}=\cos t+J\sin t,
 \qquad e^{tJ}=1+tJ,
 \qquad e^{tJ}=\cosh t+J\sinh t,
\]
y en el sector partido aparecen los proyectores
\(P_\pm=(1\pm J)/2\). Estas identidades se verifican exactamente en la
realizaci\'on declarada. La
curvatura no procede del signo aislado: aparece cuando APP aporta enlaces y
dos-celdas y se fija una conexi\'on. En la polarizaci\'on mixta suma--producto
publicada, sobre las \(81\) plaquetas, la curvatura orientada satisface
\[
 F(S,S)=F(P,P)=0,
 \qquad F(S,P)=+1,
 \qquad F(P,S)=-1\pmod9.
\]
La identidad de Stokes finita correspondiente controla que el borde de cada
rect\'angulo reproduzca la suma de sus plaquetas.

Al ordenar las transiciones se obtiene el lector temporal tr\'itico:
\[
 \begin{aligned}
 +1&\longmapsto\text{retorno, memoria y pasado},\\
  0&\longmapsto\text{umbral torsional y presente},\\
 -1&\longmapsto\text{apertura bidireccional y futuro}.
 \end{aligned}
\]
Su precursor exacto es la diferencia entre retorno de fase y retorno de
estado:
\[
 g(k+9)=g(k),
 \qquad B_{k+9}\neq B_k\quad\text{en general}.
\]
En la realizaci\'on geom\'etrica correspondiente se habla de sector
esf\'erico--volum\'etrico, umbral plano con registro torsional y sector
hiperb\'olico--areal. Esta correspondencia define un lector de Mec\'anica
Dimensional sujeto al mapa entre el TRIT y la geometr\'ia; no redefine el TRIT
algebraico ni identifica por s\'i sola
un signo con una curvatura del espacio-tiempo. La expresi\'on ``cristal de
tiempo'' nombra esta arquitectura en la terminología HMT; una realizaci\'on
din\'amica de cristal de tiempo exige adem\'as un Hamiltoniano o un protocolo
peri\'odico de excitaci\'on,
respuesta subarm\'onica robusta y estabilidad.

\subsection{Soporte exacto de la lectura ida--retorno}

El intercambio bidireccional no descansa s\'olo en una imagen verbal. Las
coordenadas angulares enteras
\[
 W_+=6n_A+s_C,
 \qquad W_-=6n_A-s_C
\]
se obtienen mediante
\[
 \binom{W_+}{W_-}
 =\begin{pmatrix}6&1\\6&-1\end{pmatrix}
 \binom{n_A}{s_C},
 \qquad
 \operatorname{SNF}\!\begin{pmatrix}6&1\\6&-1\end{pmatrix}
 =\operatorname{diag}(1,12).
\]
La ret\'icula tiene, por tanto, \`indice doce. La involuci\'on
\(s_C\mapsto-s_C\) intercambia \(W_+\leftrightarrow W_-\) y, de forma
equivalente, \(q_+\leftrightarrow q_-\). Este es el sustrato algebraico exacto
de ida y retorno. Su interpretaci\'on como propagaci\'on avanzada y retardada
requiere el operador de evoluci\'on de la representaci\'on f\'isica.

La proyecci\'on par empleada en la lectura biangular satisface, para toda
funci\'on para la que las expresiones est\'en definidas,
\[
 \frac{f(\phi)+f(-\phi)}2=f_{\rm par}(\phi),
 \qquad
 \frac{(\pi-\phi)^2+(\pi+\phi)^2}{2}=\pi^2+\phi^2.
\]
Son identidades exactas de simetrizaci\'on. La lectura en t\'erminos de una
teor\'ia f\'isica del absorbedor pertenece al funtor de propagaci\'on, no a la
identidad algebraica aislada.

\subsection{Periodicidades espinoriales inequivalentes de \texorpdfstring{\(2\pi\)}{2 pi} y \texorpdfstring{\(4\pi\)}{4 pi}}

En el levantamiento espinorial cl\'asico,
\[
 U(\theta)=\exp\!\left(-\frac{i\theta}{2}\sigma_z\right),
 \qquad U(2\pi)=-I,
 \qquad U(4\pi)=I.
\]
La igualdad es exacta en la doble cubierta \(SU(2)\to SO(3)\). Su aplicaci\'on
a una secci\'on HMT requiere especificar la representaci\'on espinorial del
transporte. Una vez fijada, explica por qu\'e una vuelta cambia el signo de la
secci\'on y dos vueltas lo restauran. Si se adoptan las relaciones CAR en el
\'algebra de ocupaci\'on, la exclusi\'on se deduce de ellas; ni CAR ni Pauli se
infieren del signo de \(2\pi\) sin esa estructura adicional.

La interpretación HMT ``hoja areal \(2\pi\), hoja volum\'etrica \(4\pi\)'' es
otro objeto. El apartado~\ref{sec:bisagra-hojas-pi-phi2} formaliza su n\'ucleo
topol\'ogico: una vuelta de la base cierra el borde proyectado en \(2\pi\) y
termina en la hoja conjugada, mientras el lazo cuadrado cierra la cubierta
orientada en \(4\pi\). El teorema es relativo al levantamiento, a la holonom\'ia y al
lector geom\'etrico declarados. Su identificaci\'on adicional con una
representaci\'on espinorial sigue requiriendo el funtor f\'isico y no se deduce
de compartir los mismos factores.

\subsection{Mapa entre multisecciones internas y representaciones físicas}

El objeto part\'icula, la proyecci\'on
extensi\'on--ruta, el selector interno de las 81 celdas, el centro electr\'onico,
las trece multisecciones finitas de familia, su elevaci\'on relativa al lector
celular \(\pi_{81}\), el tr\'itico de color, el registro, la firma, el car\'acter y
las torres forman la composición
\[
\text{extensi\'on}\longrightarrow\text{ruta}\longrightarrow
\text{multisecci\'on}\longrightarrow\text{firma}\longrightarrow
\text{car\'acter}\longrightarrow\text{torres}.
\]

La obligación restante es construir un funtor de
realizaci\'on que empareje ciertas etiquetas f\'isicas individualizadas y
compuestas con las multisecciones enriquecidas pertinentes y reproduzca las
firmas refinadas declaradas sin consultar masas ni firmas objetivo. Para los
estados compuestos, el registro se compone antes de proyectarse a cinco
coordenadas. El dominio del funtor conserva el electr\'on central, el atlas de
familias y la ley interna de car\'acter.

\enlargethispage{6\baselineskip}
\paragraph{Tipado probatorio de persistencia, atlas y familias.}
La persistencia de ruta, las anti-secciones, la lectura bidireccional y el
lector temporal tr\'itico fijan la arquitectura de partida. El sistema
inverso, el selector interno y la separaci\'on de tipos realizan esa
arquitectura mediante morfismos declarados. Los censos del atlas y las
identidades de \(C_M\) sobre palabras se cumplen en todo el sistema finito;
su levantamiento a secciones y la banda de M\"obius se demuestran con
estructura de partida expl\'icita. Los coeficientes electr\'onicos
\((-22,+15)\) se deducen con la misma estructura expl\'icita. Si una
elevaci\'on abre previamente una diferencia metrol\'ogica o una firma objetivo,
su valor es una reconstrucci\'on calibrada.

\endgroup

\Needspace{18\baselineskip}
\section{Ontología material, carga, espín, estadística, color y sabor}
La definición plena se completa con los dominios de representación, las
condiciones de conjugación, las completaciones CAR/CCR, los sectores de
calibre, la composición, la virtualidad y los controles de positividad.

\begingroup
\let\chapter\section
\let\section\subsection
\let\subsection\subsubsection
\let\subsubsection\paragraph
\let\clearpage\relax
\section{Ontología HMT--MD de las partículas}
\Needspace{7\baselineskip}
\subsection{Estado genealógico material: extensión, ruta, registro, firma y representación}
Un estado HMT completo no es un punto desnudo. Es la tupla

\[
 \mathfrak x=
 (\Gamma,\mathcal L_\Gamma,\sigma_\Gamma,\eta_\Gamma,
 \mathcal R_\Gamma,\mathcal B_\Gamma,\mathcal C_\Gamma),
\]
donde \(\Gamma\) es una ruta superviviente, \(\mathcal L_\Gamma\) su registro genealógico, \(\sigma_\Gamma\) su firma entera, \(\eta_\Gamma\) su refinamiento armónico, \(\mathcal R_\Gamma\) su representación, \(\mathcal B_\Gamma\) su memoria de retorno y \(\mathcal C_\Gamma\) su dato de composición o frontera. Dos estados con igual palabra visible pueden ser distintos si difieren en prefijo, bloque, memoria, acarreo, hoja o sombra de Witt.

\Needspace{7\baselineskip}
\subsection{Partícula como sección local persistente del estado genealógico dimensional}
Una partícula es una clase de equivalencia de secciones persistentes que satisface cuatro condiciones:

\begin{enumerate}
\item prolongación compatible a través de los cilindros del TPK;
\item retorno de fase con memoria de estado;
\item representación física irreducible o bloque mínimo invariante;
\item estabilidad espectral bajo el acoplamiento que la observa.
\end{enumerate}
Una sección persistente es una familia compatible

\[
 s=(s_K)_{K\ge K_0},\qquad
 p_{L,K}(s_L)=s_K\quad(L\ge K\ge K_0).
\]
Dos secciones \(s\) y \(s'\) representan la misma partícula cuando existe un nivel \(K_1\) y una familia compatible de simetrías internas \(g_K\in G_K\) tal que

\[
 s'_K=g_Ks_K\quad(K\ge K_1),
 \qquad
 p_{L,K}g_L=g_Kp_{L,K},
\]
y los caracteres observables del registro genealógico coinciden. Esta relación identifica descripciones que difieren por calibre interno, pero no identifica rutas con memoria, acarreo, hoja, bloque o holonomía distintos. La partícula es la clase

\[
 [s]_{\rm pers}
 =\{s':s'\sim s\},
\]
no una palabra visible aislada ni una celda elegida después de conocer su masa.

Si \(\widetilde\Gamma_K\) es la extensión al nivel \(K\), la persistencia se expresa por

\[
 p_{K+9,K}(\widetilde\Gamma_{K+9})
 =\widetilde\Gamma_K,\qquad
 g(K+9)=g(K),
\]
sin exigir

\[
 \mathcal B_{K+9}=\mathcal B_K.
\]
La partícula no es la fase que retorna, sino la identidad genealógica conservada a través del cambio de estado. Su masa mide la rigidez espectral de esa persistencia cuando se acopla a una realización material.

\Needspace{7\baselineskip}
\subsection{Antipartícula como imagen conjugada de una sección persistente bajo inversión de hoja}
La conjugación dodecafásica es

\[
 C_M(\mathbf t)=-\operatorname{rev}(\mathbf t).
\]
El par partícula--antipartícula está formado por dos orientaciones de una misma banda genealógica. Si \(\chi\) denota la orientación de hoja,

\[
 (\mathbf t,\chi)\sim(C_M\mathbf t,-\chi).
\]
Para una energía de banda

\[
 E_{\rm band}(\mathbf t,\chi)
 =E_+(\mathbf t)+\chi E_-(\mathbf t),
\]
con

\[
 E_\pm(\mathbf t)
 =\frac12\bigl(E_0(\mathbf t)\pm E_0(C_M\mathbf t)\bigr),
\]
se obtiene

\[
 E_{\rm band}(C_M\mathbf t,-\chi)
 =E_{\rm band}(\mathbf t,\chi).
\]
La igualdad de masas es, por tanto, consecuencia de la paridad del operador bajo la conjugación, mientras que carga y orientación cambian. La igualdad de anchuras requiere además que el operador de decaimiento sea conjugado por \(C_M\).

\Needspace{7\baselineskip}
\subsection{Masa como carácter espectral de persistencia de \(1\) ruta material}
La masa es la medida espectral, en la recta \(\mathcal L_M\), del operador que resulta de comprimir carácter, torres, lector bidireccional y reloj sobre el subespacio persistente seleccionado por el acoplamiento:

\[
 \widehat M_{P,a}
 =P_{P,a}\,
 \mathbf m_\star^{\rm ret}
 R_{\rm act}^{\widehat K_{P,a}/2}
 \widehat{\mathcal A}_{P,a}^{(0)}
 R_{\rm act}^{\widehat K_{P,a}/2}
 P_{P,a}.
\]
La salida general es

\[
 \mu_{P,a}^{\rm mass}(B)
 =\operatorname{Tr}\!\left(
 \rho_{P,a}E_{\widehat M_{P,a}}(B)\right).
\]
Existe una masa escalar \(m_P\) si y sólo si

\[
 P_{\operatorname{supp}\rho}\widehat M_{P,a}
 P_{\operatorname{supp}\rho}
 =m_P P_{\operatorname{supp}\rho},
\]
equivalentemente, si la varianza de \(\widehat M_{P,a}\) en el estado preparado es cero. Esta definición incluye el electrón de rango uno, los multipletes forzados por Schur y los polos aislados de resonancias, sin confundirlos.

\Needspace{7\baselineskip}
\subsection{Carga como carácter de calibre transportado por la sección material}
La carga es el carácter orientado de la capa impar del registro genealógico bajo conjugación. Sea \(\sigma_{\rm odd}\) la firma impar. Entonces

\[
 Q(C_M\Gamma)=-Q(\Gamma),\qquad
 Q(\Gamma)=\mathfrak q\!\left(\sigma_{\rm odd}(\Gamma)\right),
\]
donde \(\mathfrak q\) es el homomorfismo de la firma impar al retículo de carga de la realización. La firma másica par y la firma de carga impar son objetos distintos: una palabra puede tener suma lineal nula y, sin embargo, portar una firma másica no nula.

\Needspace{7\baselineskip}
\subsection{Espín como representación de la acción rotacional sobre la fibra material}
El espín es la clase de holonomía con la que una sección persistente responde a una rotación completa del marco de orientación. La raíz interna

\[
 J^2=-I
\]
separa los dos cierres:

\[
 U(2\pi)=
 \begin{cases}
 +I,&\text{sector tensorial},\\
 -I,&\text{sector espinorial},
 \end{cases}
 \qquad
 U(4\pi)=I.
\]
Un bosón entero vive en una representación que cierra a \(2\pi\); un fermión semientero necesita la doble vuelta \(4\pi\). El valor de espín no se deduce de un conteo cardinal aislado: es el peso de la representación del grupo de holonomía sobre el subespacio persistente. En particular:

\[
 s=0:\text{ sección escalar},\qquad
 s=\tfrac12:\text{ banda espinorial},
\]
\[
 s=1:\text{ sección vectorial transversal},\qquad
 s=2:\text{ componente tensorial de una realización métrica}.
\]
El último renglón define un tipo de representación; no obliga a que exista una partícula elemental de espín dos.

La quiralidad y la helicidad son datos posteriores y distintos. Una involución orientada \(\Gamma_{\rm ch}^2=I\) define

\[
 P_L=\frac{I-\Gamma_{\rm ch}}2,
 \qquad
 P_R=\frac{I+\Gamma_{\rm ch}}2,
\]
mientras que la helicidad es el signo espectral del generador de rotación sobre la dirección de propagación. La primera clasifica hojas de la representación; la segunda depende del estado cinemático. Esta separación permite que un neutrino sea fermiónico por su cierre \(4\pi\), aunque su acoplamiento débil seleccione una sola quiralidad.

\Needspace{7\baselineskip}
\subsection{Estadística cuántica inducida por la simetría de intercambio de secciones}
La estadística es la regla de completación del espacio de rutas idénticas:

\[
 \mathcal F_-(\mathcal H)
 =\bigoplus_{n\ge0}\wedge^n\mathcal H,\qquad
 \mathcal F_+(\mathcal H)
 =\bigoplus_{n\ge0}\operatorname{Sym}^n\mathcal H.
\]
En la completación exterior, las relaciones CAR

\[
 \{a_i,a_j^\dagger\}=\delta_{ij},\qquad
 \{a_i,a_j\}=0
\]
implican

\[
 n_i^2=n_i
\]
y, por tanto, exclusión de Pauli y distribución de Fermi--Dirac. En la completación simétrica, las relaciones CCR permiten ocupación múltiple y conducen a Bose--Einstein. La estadística no se infiere de la masa: procede de cómo se compone la ruta bajo intercambio.

\Needspace{7\baselineskip}
\subsection{Color como representación residual de \(SU(3)\) sobre las rutas quark}
El color es la órbita residual triádica de una sección quark dentro de la fibra de carga y sabor. Las treinta y seis posiciones coloreadas se descomponen como

\[
 12_R+12_G+12_B.
\]
El objeto físico no es uno de esos colores, sino el triplete y su acción:

\[
 \mathcal H_q^{\rm color}
 =\mathcal H_R\oplus\mathcal H_G\oplus\mathcal H_B.
\]
Si el operador de masa conmuta con la acción irreducible de color,

\[
 U_g\widehat M_qU_g^*=\widehat M_q,
\]
el lema de Schur fuerza una masa común en el triplete. Elegir una celda roja, verde o azul como representante destruiría la covariancia.

\Needspace{7\baselineskip}
\subsection{Sabor como clase de ruta dentro de una generación material}
El sabor es una clase de acoplamiento inequivalente dentro de una misma representación de calibre, caracterizada por carga, generación, hoja, firma fina, memoria y proyector de interacción. Dos sabores pueden compartir una fibra gruesa y diferir en el entrelazador que los acopla al lector.

Formalmente, si \(F\) es la fibra común, los sabores \(f_i\) son representaciones \(\pi_i\) para las que

\[
 \operatorname{Hom}_{G_F}(V_{f_i},\ell^2(F))\ne0
\]
y cuyos proyectores \(P_{f_i}\) no son conjugados por una simetría de calibre que se mantenga física. El sabor no es la profundidad del atlas y no se reduce a un nombre convencional.

\Needspace{7\baselineskip}
\subsection{Generación material como estrato recurrente tipado por profundidad de retorno, firma fina y espectro del lector}
Una generación es un estrato recurrente de la misma gramática de carga y representación, distinguido por profundidad de retorno, firma fina y espectro del lector. La etiqueta de generación sólo es física cuando existe un proyector que entrelaza ese estrato con el canal de acoplamiento. Por ello, una profundidad \(G_j\) no debe identificarse automáticamente con electrón, muón, tau o neutrino estéril.

\Needspace{7\baselineskip}
\subsection{Campo de calibre como conexión sobre las fibras de carga}
Un campo de calibre es la realización de una redundancia de marco sobre las fibras. Sus bosones elementales pertenecen a la dirección nula del operador de masa mientras la simetría permanezca exacta. Una masa de calibre no se obtiene aproximando a cero un carácter positivo, sino mediante una ruptura, un acoplamiento longitudinal o una realización de polo declarada.

\Needspace{7\baselineskip}
\subsection{Partícula compuesta como sección persistente del producto de rutas enlazadas}
Una partícula compuesta es una sección persistente del operador de composición, no la suma de las masas ni de las firmas de sus constituyentes. Si \(\Gamma_1,\ldots,\Gamma_n\) son rutas constituyentes, primero se construye

\[
 \mathfrak C_{12}
 (\Gamma_1,\ldots,\Gamma_n;\mathcal T,\partial,\mathrm{carry})
 =\mathcal L_{\rm comp},
\]
donde \(\mathcal T\) es el árbol temporal y \(\partial\) la frontera de enlace. Sólo después actúa el lector

\[
 \mathcal L_{\rm comp}\longrightarrow\sigma_{\rm comp}
 \longrightarrow\chi_{\rm comp}
 \longrightarrow\widehat M_{\rm comp}.
\]
La energía de ligadura reside en el cociclo de composición, la memoria de frontera y la modificación del espectro; no se introduce como corrección externa.

\Needspace{7\baselineskip}
\subsection{Resonancia como sección transitoria y anchura como tasa de desintegración}
Una resonancia es un modo persistente no unitario del operador acoplado. Si

\[
 \lambda=re^{-i\theta},\qquad 0<r<1,
\]
entonces

\[
 E=\frac{\hbar_{\rm int}}{t_0}
 \bigl(\theta-i\log r\bigr)
 =M-\frac{i}{2}\Gamma.
\]
La masa \(M\) y la anchura \(\Gamma\) son las partes real e imaginaria de un polo, pero requieren lectores distintos. La anchura también puede obtenerse de un semigrupo comprimido:

\[
 T_X=Q_XUQ_X,\qquad
 \lambda_X=-\frac1{t_0}\log\rho(T_X),\qquad
 \Gamma_X=\hbar_{\rm int}\lambda_X.
\]
No se deduce una anchura únicamente de que una ruta posea horizonte finito.

\Needspace{7\baselineskip}
\subsection{Sección virtual como extensión intermedia no asintótica del atlas}
Una sección es virtual cuando su prolongación termina de manera no trivial:

\[
 \operatorname{Ext}_{L}(s)\ne\varnothing,\qquad
 \operatorname{Ext}_{L+1}(s)=\varnothing,\qquad
 \mathcal L(s)=L<\infty.
\]
Su contenido primario es el horizonte \(L\), la obstrucción terminal, el registro genealógico finito y la amplitud de propagación. No es una partícula de masa imaginaria ni un estado documental todavía no clasificado.

\Needspace{7\baselineskip}

\endgroup
\begingroup
\let\chapter\section
\let\section\subsection
\let\subsection\subsubsection
\let\subsubsection\paragraph
\let\clearpage\relax
\section{Clasificación completa de partículas y estados físicos}
\label{sec:catalogo-realizaciones-tipadas-rev8}
\index{sectores físicos!inventario tipado}
\index{atlas!81/56/13/324}

Las realizaciones se ordenan por objeto, multiplicidad y codominio. El dominio matemático contiene \(81\) posiciones, agrupadas en
codominios distintos. El dominio matemático contiene \(81\) posiciones, agrupadas en
\(56\) familias de trayectoria y \(13\) clases de extensión; las \(324\)
trayectorias orientadas añaden la orientación inicial y no representan masas
ni partículas adicionales. La columna ``operador de evaluación'' no impone una masa donde el objeto exige
una carta nula, una composición previa, un invariante topológico o una
obstrucción de prolongación.

\subsection{Definiciones espectrales de partícula, espín, color y sabor}

Una \emph{partícula} es una sección superviviente realizada como ruta
orientada, dotada de registro, representación de simetría, operador espectral
y carta física. Un fermión realiza la graduación exterior y las relaciones
CAR; un bosón realiza la completación simétrica y las relaciones CCR. El
espín registra la holonomía de la orientación: el retorno de \(2\pi\) cierra
la fase bosónica, mientras la sección espinorial recupera su estado tras
\(4\pi\). Esta distinción enlaza la doble hoja areal--volumétrica con la
paridad y la exclusión de Pauli.

El \emph{color} es la representación ternaria que eleva la coordenada
\(r-c\pmod3\) a la órbita fundamental y, después, al adjunto de
\(\mathfrak{su}_3\). El \emph{sabor} identifica una base de realización
dentro de una representación con carga, color y espín fijados; CKM y PMNS
transportan entre la base de interacción y la base espectral. Los neutrinos
son secciones fermiónicas neutras: sus etiquetas de sabor pertenecen al
lector de mezcla y sus masas a los autovectores del operador correspondiente.

Una partícula virtual es una sección finita con registro y propagador hasta
un horizonte de prolongación; una realización Majorana es el sector
fermiónico real fijado por conjugación; y los sectores de Fibonacci, Ising y
\(\mathbb Z/6\mathbb Z\) publican dimensiones cuánticas, reglas de fusión y
caracteres de trenza. Cada definición conserva así el codominio que la
estructura le asigna.


\endgroup
\begingroup
\let\chapter\section
\let\section\subsection
\let\subsection\subsubsection
\let\subsubsection\paragraph
\let\clearpage\relax

La tabla conserva además las realizaciones cuánticas que su clasificación
abrevia. Los tres leptones cargados y los seis sabores quark pertenecen a la
completación exterior: CAR conduce a Pauli y a Fermi--Dirac, sin atribuir esas
relaciones a un TRIT aislado. Los quarks ocupan treinta y seis celdas al
conservar color y antipartícula antes del cociente. El soporte neutro común
\(\mathcal H_\nu=\bigoplus_{j=1}^{4}\mathcal H_{\nu,G_j}\) tiene dimensiones
\(2+4+6+8=20\). Los símbolos \(G_j\) designan estratos de profundidad del atlas;
\(f_e,f_\mu,f_\tau\) designan la base de interacción. En particular, \(G_4\) es
un estrato del soporte neutro común y no se identifica con un estado estéril o
exótico sin un proyector adicional. PMNS transporta la base de interacción a la
base de autovectores; no crea autovalores, escala de masa ni selector de
profundidad. El fotón conserva ida y retorno dentro de la
carta gauge nula: \(m=0\) no significa evaluar \(\chi_0\) con un exponente muy
negativo. \(W^\pm\) y \(Z^0\) son representaciones bosónicas masivas, con
conjugación de carga para \(W^+\leftrightarrow W^-\). Por último, la
antisección se obtiene mediante inversión y reversión de banda; la masa es par
bajo esa involución y no se duplica como un segundo parámetro independiente.

Todos los sectores conservan el prefijo causal
\[
\text{extensión superviviente}\longrightarrow\text{ruta observable}
\longrightarrow\text{registro cronológico enriquecido}.
\]
Después cambia el lector. En el sector material positivo continúa por
\(\text{firma}\to\text{carácter}\to\text{torres}\to\text{carta dimensional}\).
El sector nulo evita el carácter positivo y fija \(m=0\); los hadrones componen
primero registros cronológicos enriquecidos; los sectores topológicos usan invariantes de
fusión o trenza; y la virtualidad se reconoce por obstrucción terminal. El
dominio de estos mapas es la clasificación completa anterior. Un ensayo
numérico sobre un subconjunto sólo establece el caso finito ensayado y no
reduce ese dominio.

Los censos que acompañan esta tabla pertenecen a cocientes diferentes:
\[
81=25_{\ell^-}+20_\nu+20_d+16_u,
\]
\[
56\ \text{familias con fibras}\quad
41\times1+10\times2+2\times3+1\times4+2\times5=81,
\]
\[
13\ \text{multisecciones},\qquad324=81\times4\ \text{genealogías}.
\]
Asimismo, las \(192\)-A clases de representación asociadas a
\(\operatorname{diag}(2,6,2,2,4)\) y el índice reticular \(192\)-B de
\(\operatorname{SNF}(2,2,2,2,12)\), con \(52\) firmas crudas, no son el
El cardinal del catálogo expositivo histórico no cuenta partículas individuales ni define la taxonomía material.
no partículas individuales.

\endgroup

\begingroup
\let\chapter\section
\let\section\subsection
\let\subsection\subsubsection
\let\subsubsection\paragraph
\let\clearpage\relax
% Vista científica exacta materializada desde una rama condicional.
% Fuente: source/demostraciones_primarias/27_04_particula_estadistica.tex; rama: HMTIncludeParticleCore.
% No modifica el contenido matemático de la rama de origen.
\chapter{Secciones persistentes y propagación finita}
\label{chap:particula}

La noción matemática de partícula se introduce después de la estructura de
frontera y antes de especializar su dinámica. Su soporte primario es una
extensión superviviente del sistema de estados TPK\@. Una trayectoria
orientada es la sombra secuencial de esa extensión; sobre ella, los grados
internos forman un fibrado discreto. El propagador, la suma de caminos y las
representaciones de Clifford se incorporan después de construir la recta de
acción. Esta jerarquía impide que una partícula se defina mediante la elección
externa de una fila o de una ruta y separa el objeto matemático de su posterior
identificación con una especie física.

\section{Extensión superviviente y proyección observable}

Sea \((\mathcal S_m,\rho_{m+1,m})\) el sistema inverso de estados que
superan los cierres hasta profundidad \(m\). Un estado global es una familia
compatible
\[
 x=(x_m)_m\in\Omega_{\HMT}^{\rm surv}
 :=\varprojlim_m\mathcal S_m.
\]
Una carta secuencial de profundidad \(R\) produce una ruta
\[
 \gamma_x^{[R]}=\operatorname{sh}_R(x)
 =(c_0\longrightarrow\cdots\longrightarrow c_R).
\]
El símbolo \(\operatorname{sh}\) denota una lectura de sombra: diferentes
cartas pueden presentar un mismo estado global, mientras conserven hoja,
orientación, borde y registro. En el
\cref{chap:extension-ruta-caracter} se construye el registro enriquecido que
hace descender el carácter de acción desde estas presentaciones a la
extensión.

\section{Partícula como sección compatible de un fibrado discreto}

Sea \(\Gamma_{\rm TPK}\) el grafo cuyos vértices son estados completos del
núcleo de fase topológica y cuyas aristas son transformaciones admisibles.
Para \(x\in\Omega_{\HMT}^{\rm surv}\) y su trayectoria observada
\[
 \gamma_x=(c_0\longrightarrow c_1\longrightarrow\cdots
 \longrightarrow c_R),
\]
sea \(\operatorname{Typ}(c_t)\) el conjunto finito de tipos aritméticos,
orientaciones y memorias sobre el estado \(c_t\), y sea
\(E_{c_t}=\mathbb C[\operatorname{Typ}(c_t)]\) el espacio vectorial complejo
libre generado por esos tipos. El fibrado discreto restringido a \(\gamma_x\) es
\[
 \mathcal B_{\gamma_x}=\bigsqcup_{t=0}^{R}E_{c_t}
 \longrightarrow\{0,\ldots,R\}.
\]
Cada arista \(c_t\to c_{t+1}\) induce una aplicación lineal de transporte
\(\tau_t:E_{c_t}\to E_{c_{t+1}}\).

\begin{definition}[Partícula matemática]
Una partícula de Mecánica Dimensional es una tupla
\[
 \mathfrak p=(x,\gamma_x,\mathcal B_{\gamma_x},s,n,U),
\]
donde \(x\in\Omega_{\HMT}^{\rm surv}\) es la extensión superviviente y
\(\gamma_x=\operatorname{sh}_R(x)\) su sombra en la carta declarada; además,
\[
 s:\{0,\ldots,R\}\longrightarrow\mathcal B_{\gamma_x},\qquad
 s(t)\in E_{c_t},\qquad s(t+1)=\tau_t(s(t)),
\]
es una sección compatible;
\[
 n:\bigsqcup_{t=0}^{R}\operatorname{Modes}(E_{c_t})
 \longrightarrow\mathbb N
\]
es su función de ocupación, donde
\(\operatorname{Modes}(E_{c_t}):=\operatorname{Typ}(c_t)\) denota la base
distinguida del espacio vectorial libre; y
\[
 U_{t_2,t_1}:E_{c_{t_1}}\longrightarrow E_{c_{t_2}}
\]
es una familia de propagadores lineales que entrelaza los transportes,
\(U_{t,t}=\operatorname{id}_{E_{c_t}}\),
\(U_{t+1,t}=\tau_t\), y satisface
\[
 U_{t_3,t_2}\circ U_{t_2,t_1}=U_{t_3,t_1}
 \qquad(0\leq t_1\leq t_2\leq t_3\leq R).
\]
Cada composición está definida entre las fibras correspondientes.
La identificación de \(\mathfrak p\) con una especie
física requiere, además, una representación del grupo de simetría pertinente
y un mapa desde estados HMT hacia los estados de dicha representación.
\end{definition}

La extensión conserva la genealogía y el registro; la ruta expone su lectura
ordenada; los grados internos pertenecen a las fibras, la ocupación a \(n\) y
la evolución a \(U\). La partícula es la compatibilidad de estos niveles, no
uno de ellos por separado. En particular, la ruta no es escogida por la
partícula: es una proyección de la extensión que la realiza.

\endgroup
\begingroup
\let\chapter\section
\let\section\subsection
\let\subsection\subsubsection
\let\subsubsection\paragraph
\let\clearpage\relax
% Fragmento científico exclusivo de 01_electron.tex, líneas 120--197.
% El prefijo 1--118 coincide byte a byte con la ontología principal ya
% publicada; este fragmento conserva únicamente la formulación cuántica
% adicional sin repetir definiciones idénticas.
\section{Ontología cuántica de la partícula}
\Needspace{7\baselineskip}
\subsection{Estado, historia y realización}
Un estado genealógico es una sección compatible del sistema de extensiones del TPK. Una historia observable es la proyección cronológica de esa sección sobre un lector. Una realización física es la representación de la historia en una recta dimensional, un espacio de Hilbert, un álgebra de campos o un codominio topológico. Estas tres nociones no se intercambian. El estado porta la identidad; la historia porta la evidencia; la realización porta el observable.

Sea \(\mathcal X_{\rm TPK}^{\rm enr}\) el espacio de estados enriquecidos y \(p\) el lector observable. Una partícula es una clase \([X]\) de secciones persistentes tal que

\[
p(X_{K+1})|_K=p(X_K),
\]
la ley de transición conserva las relaciones del registro genealógico y la prolongación no termina en un horizonte de obstrucción. La ruta

\[
\gamma_X=(p(X_0),p(X_1),p(X_2),\ldots)
\]
es la historia visible de la partícula. No es una trayectoria espacial por definición. Una realización métrica posterior puede convertirla en curva, órbita o propagador cuando construye el entrelazador correspondiente.

La identidad física se conserva bajo refinamiento cuando los cuadrados

\[
\begin{CD}
X_{K+1} @>{F_{K+1}}>> X'_{K+1}\\
@V{r_{K+1,K}}VV @VV{r'_{K+1,K}}V\\
X_K @>>{F_K}> X'_K
\end{CD}
\]
conmutan. La masa, la carga y el espín son lectores diferentes de ese mismo antecedente. La igualdad de uno no obliga a la igualdad de los demás.

\Needspace{7\baselineskip}
\subsection{Representación cuántica de la partícula como sección local persistente del estado genealógico}
Una partícula HMT--MD es una sección persistente de una extensión enriquecida que satisface cuatro condiciones:

\begin{enumerate}
\item posee una ruta observable compatible con la monodromía nonádica;
\item conserva un registro genealógico bajo refinamiento;
\item admite una representación de orientación, carga, espín y estadística;
\item presenta una salida espectral estable en el codominio físico declarado.
\end{enumerate}
La estabilidad no significa inmovilidad. Significa que los cambios de fase, hoja y memoria pertenecen a una misma clase genealógica. La partícula puede cerrar su fase y no su estado, cambiar de hoja y conservar su masa, o atravesar varias celdas y mantener la misma identidad de extensión.

Una especie es una familia de partículas que comparte el clasificador interno

\[
\mathfrak s=
(\mathfrak q,\mathfrak j,\mathfrak c,\mathfrak f,\mathfrak g,
\mathfrak r,\mathfrak o),
\]
donde \(\mathfrak q\) es la representación de carga, \(\mathfrak j\) la holonomía de espín, \(\mathfrak c\) la órbita de color, \(\mathfrak f\) el sabor, \(\mathfrak g\) la generación, \(\mathfrak r\) la familia de ruta y \(\mathfrak o\) el operador de masa. Una especie no central puede corresponder a una multisección completa; exigirle un único punto APP antes de declarar el acoplamiento físico constituye un error de tipo.

\Needspace{7\baselineskip}
\subsection{Conjugación cuántica de la antipartícula mediante inversión de hoja de una sección persistente}
La involución dodecafásica es

\[
C_M(\tau)=-\operatorname{rev}(\tau).
\]
Actúa sobre la historia orientada y se prolonga al registro genealógico. La carga y la orientación lineal son impares; la memoria cuadrática es par. Si \(q\) es el lector de carga,

\[
q(C_MX)=-q(X).
\]
La igualdad de masas exige además covariancia operatoria. Si \(K\) realiza la conjugación en el espacio físico y

\[
K\widehat M_XK^{-1}=\widehat M_{\bar X},
\]
transportando también los canales y la prescripción causal cuando \(K\) es antiunitario, entonces \(m_{\bar X}=m_X\). La paridad del registro por sí sola no sustituye esta premisa.

Partícula y antipartícula son dos secciones orientadas de una misma banda genealógica. En la realización de Möbius, una vuelta intercambia la hoja y dos vueltas restauran la orientación. La conclusión no procede de dibujar una cinta torcida: exige la involución, el levantamiento a la doble cubierta, el lazo de retorno y la holonomía de signo.

Para el registro electrónico \(\tau_e=+++---+++---\),

\[
C_M\tau_e=\tau_e.
\]
La palabra central es fija, mientras la representación de carga distingue electrón y positrón. Así se obtiene una misma masa sin duplicar la ruta temporal. El cierre a 216 pasos no significa 108 pasos de partícula y 108 de antipartícula: significa restauración del levantamiento cinemático orientado.

\Needspace{7\baselineskip}

\endgroup
\begingroup
\let\chapter\section
\let\section\subsection
\let\subsection\subsubsection
\let\subsubsection\paragraph
\let\clearpage\relax
La Ley General de Masas queda formulada como una transformación genealógica y
espectral, no como una colección de números. Su dominio efectivo es la cadena
\[
\text{APP}\longrightarrow\text{TRIT}\longrightarrow\text{TPK}
\longrightarrow\Gamma^{\rm surv}\longrightarrow\Gamma^{\rm obs}
\longrightarrow L_{108}\longrightarrow\sigma\in\mathbb Z^5
\longrightarrow\chi_{\rm full}\longrightarrow\widehat M,
\]
seguida, cuando la representación física lo permite, por el funcional que
extrae una coordenada escalar de la recta de masa. Cada flecha conserva su
dominio, su orientación, su memoria y su estatuto; ninguna cifra externa elige
retroactivamente una ruta, una firma o un coeficiente armónico.

El atlas finito contiene 81 celdas, 56 familias de ruta, 13 multisecciones y
324 rutas orientadas. El punto central electrónico es la única fibra puntual
equivariante. Las demás clases elementales se realizan naturalmente como
fibras, órbitas o multisecciones, de modo que su salida primaria puede ser un
operador o un espectro antes de ser una masa escalar. Esta diferencia no es una
carencia de dominio, sino una propiedad de la geometría interna.

En la fibra central, el selector basal y el lector bidireccional producen
\[
\mathfrak e_0
=\frac{\sqrt3}{4}
 \left[\exp\!\left(\frac{\varphi}{\pi^2}\right)
       -22\alpha_{\rm HMT}^3\right](1+15\Delta_4),
\qquad
\kappa_e=80-54\alpha_{\rm HMT}+6\Delta_4,
\]
\[
\mathfrak e_e^\beta=\mathfrak e_0R_{\rm act}^{\kappa_e},
\qquad
\varsigma_{u_E}(\mathfrak e_e^\beta)
=0.510998950690000808608\ldots\ \mathrm{MeV}.
\]
La escala absoluta de acción permanece en la rama
\(\Sigma_H=\varphi\alpha_{\rm HMT}^{16}\), de orden decimal \(-34\), y se
acopla al reloj mediante \(\hbar_{\rm int}\omega\); al formar cocientes de
masas, toda sección común cancela y sólo sobreviven los caracteres, los
lectores y los refinamientos relativos correctamente tipados.

La jerarquía armónica no consta de ocho correcciones aisladas. Es el semigrupo
global
\[
\langle90,120\rangle
=\{90,120\}\cup\{30r:r\ge6\},
\qquad
c_n=q_-^n+q_+^n,
\quad s_n=q_-^n-q_+^n,
\]
actuando sobre un único operador bicapa. Los niveles visibles son testigos de
esa jerarquía; su completitud representacional no sustituye al mapa físico que
debe transportar ruta, acarreo, memoria y devanado a los coeficientes de cada
estado.

Las realizaciones físicas tipadas, las diecisiete evaluaciones condicionadas, las
dos secciones nulas y los sectores neutrínicos, compuestos, resonantes,
topológicos y virtuales quedan así reunidos sin aplanar sus codominios. El
inventario posterior de 471 masas y 384 anchuras conserva el contraste
estado por estado, pero no se confunde con un censo de predicciones. Una
comparación cuantitativa adquiere fuerza física sólo después de declarar
observable, unidad, esquema, escala, incertidumbre, covarianza y regla de
realización.

La conclusión estructural es precisa: la masa es la medida espectral de una
persistencia bajo acoplamiento. La partícula es la sección persistente; la ruta,
su historia observable; la firma, su lectura entera; el carácter, su acción
adimensional; las torres, su resolución armónica; y la carta dimensional, la
realización de esa genealogía en unidades físicas. El atlas y los inventarios
precedentes permiten reproducir cada afirmación sin reducir la teoría a una
tabla abreviada ni convertir una interfaz abierta en un teorema cerrado.

\endgroup

\section{Composición e identidad}
Una partícula compuesta conserva identidad cuando el registro de enlace, la
frontera, la representación y los observables permanecen compatibles bajo
refinamiento. Una resonancia añade un polo y una anchura; no constituye una
nueva coordenada de la firma de masa. El criterio de identidad queda refutado
si depende del valor externo o si dos historias inequivalentes se identifican
sin un isomorfismo de secciones.
